A \code{ThreadSet} is a set of Thread pointers. It has similar
operations as Thread, and operations done on a ThreadSet affect every Thread in
that ThreadSet. One some system, such as Blue Gene Q, using a ThreadSet is
more efficient when doing the same operation across a large number of Threads.
It satisfies the C++
\href{https://en.cppreference.com/cpp/named_req/ForwardIterator}{LegacyForwardIterator}
concept. It also offers the standard \code{begin/end} members for iteration.

\begin{apient}
ThreadSet::ptr
ThreadSet::const_ptr
\end{apient}

\apidesc{
The ptr and const\_ptr types are smart pointers to a ThreadSet object. When
the last smart pointer to the ThreadSet is cleaned, then the underlying
ThreadSet is cleaned. The const\_ptr type is a const smart pointer.
}

\begin{apient}
ThreadSet::weak_ptr
ThreadSet::const_weak_ptr
\end{apient}

\apidesc{

The weak\_ptr and const\_weak\_ptr are weak smart pointers to a ThreadSet
object. Unlike regular smart pointers, weak pointers are not counted as
references when determining whether to clean the ThreadSet object. The
const\_weak\_ptr type is a const weak smart pointer.
}

\begin{apient}
static ThreadSet::ptr newThreadSet()
\end{apient}

\apidesc{
This function creates a new ThreadSet that is empty.
}

\begin{apient}
static ThreadSet::ptr newThreadSet(Thread::ptr thr)
\end{apient}

\apidesc{
This function creates a new ThreadSet that contains thr.
}

\begin{apient}
static ThreadSet::ptr newThreadSet(const ThreadPool &threadp)
\end{apient}

\apidesc{
This function creates a new ThreadSet that contains all of the Threads currently
in threadp.
}

\begin{apient}
static ThreadSet::ptr newThreadSet(const std::set<Thread::const_ptr> &thrds)
\end{apient}

\apidesc{
This function creates a new ThreadSet that contains all of the threads in thrds.
}

\begin{apient}
static ThreadSet::ptr newThreadSet(ProcessSet::ptr pset)
\end{apient}

\apidesc{
This function creates a new ThreadSet that contains every live thread currently
in every process in pset.
}

\begin{apient}
ThreadSet::ptr set_union(ThreadSet::ptr tset) const
\end{apient}

\apidesc{
This function returns a new ThreadSet whose elements are a set union of this
ThreadSet and tset.
}

\begin{apient}
ThreadSet::ptr set_intersection(ThreadSet::ptr tset) const
\end{apient}

\apidesc{
This function returns a new ThreadSet whose elements are a set intersection of
this ThreadSet and tset.
}

\begin{apient}
ThreadSet::ptr set_difference(ThreadSet::ptr tset) const
\end{apient}

\apidesc{
This function returns a new ThreadSet whose elements are a set difference of
this ThreadSet minus tset.
}

\begin{apient}
iterator find(Thread::const_ptr thr)
const_iterator find(Thread::const_ptr thr) const
\end{apient}

\apidesc{
These functions search a ThreadSet for thr and returns an iterator pointing to
that element. It returns ThreadSet::end() if no element is found
}

\begin{apient}
bool empty() const
\end{apient}

\apidesc{
This function returns true if the ThreadSet has zero elements and false
otherwise.
}

\begin{apient}
size_t size() const
\end{apient}

\apidesc{
This function returns the number of elements in the ThreadSet.
}

\begin{apient}
std::pair<iterator, bool> insert(Thread::const_ptr thr)
\end{apient}

\apidesc{
This function inserts thr into the ThreadSet. If thr already exists in the
ThreadSet, then no change will occur. This function returns an iterator
pointing to either the new or existing element and a boolean that is true if an
insert happened and false otherwise.
}

\begin{apient}
void erase(iterator pos)
\end{apient}

\apidesc{
This function removes the element pointed to by pos from the ThreadSet.
}

\begin{apient}
size_t erase(Thread::const_ptr thr)
\end{apient}

\apidesc{
This function searches the ThreadSet for thr, then erases that element from the
ThreadSet. It returns 1 if it erased an element and 0 otherwise.
}

\begin{apient}
void clear()
\end{apient}

\apidesc{
This function erases all elements in the ThreadSet.
}

\begin{apient}
ThreadSet::ptr getErrorSubset() const
\end{apient}

\apidesc{
This function returns a new ThreadSet containing every Thread from this
ThreadSet that has a non-zero error code. Error codes are reset upon every
ThreadSet API call, so this function shows which Threads had an error on the
last ThreadSet operation.
}

\begin{apient}
void getErrorSubsets(std::map<ProcControlAPI::err_t, ThreadSet::ptr> &err) const
\end{apient}

\apidesc{
This function returns a set of new ThreadSets containing every Thread from this
ThreadSet that has non-zero error codes, and grouped by error code. For each
error code generated by the last ThreadSet API operation an element will be
added to err, and every Thread that has that error code will be added to the
new ThreadSet associated with that error code.
}

\begin{apient}
bool allStopped() const
bool anyStopped() const
\end{apient}

\apidesc{
These functions respectively return true if any or all threads in this ThreadSet
are stopped and false otherwise.
}

\begin{apient}
bool allRunning() const
bool anyRunning() const
\end{apient}

\apidesc{
These functions respectively return true if any or all threads in this ThreadSet
are running and false otherwise.
}

\begin{apient}
bool allTerminated() const
bool anyTerminated() const
\end{apient}

\apidesc{
These functions respectively return true if any or all threads in this ThreadSet
are terminated and false otherwise.
}

\begin{apient}
bool allSingleStepMode() const
bool anySingleStepMode() const
\end{apient}

\apidesc{
These functions respectively return true if any or all threads in this ThreadSet
are running in single step mode, and false otherwise.
}

\begin{apient}
bool allHaveUserThreadInfo() const
bool anyHaveUserThreadInfo() const
\end{apient}

\apidesc{
These functions respectively return true if any or all threads in this ThreadSet
have user thread information available and false otherwise.
}

\begin{apient}
ThreadSet::ptr getStoppedSubset() const
\end{apient}

\apidesc{
This function returns a new ThreadSet, which is a subset of this ThreadSet, and
contains every Thread that is stopped.
}

\begin{apient}
ThreadSet::ptr getRunningSubset() const
\end{apient}

\apidesc{
This function returns a new ThreadSet, which is a subset of this ThreadSet, and
contains every Thread that is running.
}

\begin{apient}
ThreadSet::ptr getTerminatedSubset() const
\end{apient}

\apidesc{
This function returns a new ThreadSet, which is a subset of this ThreadSet, and
contains every Thread that is terminated.
}

\begin{apient}
ThreadSet::ptr getSingleStepSubset() const
\end{apient}

\apidesc{
This function returns a new ThreadSet, which is a subset of this ThreadSet, and
contains every Thread that is in single step mode.
}

\begin{apient}
ThreadSet::ptr getHaveUserThreadInfoSubset() const
\end{apient}

\apidesc{
This function returns a new ThreadSet, which is a subset of this ThreadSet, and
contains every Thread that has user thread information available.
}

\begin{apient}
bool getStartFunctions(AddressSet::ptr result) const
\end{apient}

\apidesc{
This function fills in the AddressSet pointed to by result with the address of
every start function of each Thread in this ThreadSet. This information is
only available on threads that have user thread information available.
This function return true if it succeeded for every Thread, and false otherwise.
}

\begin{apient}
bool getStackBases(AddressSet::ptr result) const
\end{apient}

\apidesc{
This function fills in the AddressSet pointed to by result with the address of
every stack base of each Thread in this ThreadSet. This information is only
available on threads that have user thread information available.
This function return true if it succeeded for every Thread, and false otherwise.
}

\begin{apient}
bool getTLSs(AddressSet::ptr result) const
\end{apient}

\apidesc{
This function fills in the AddressSet pointed to by result with the address of
every thread-local storage region of each Thread in this ThreadSet. This
information is only available on threads that have user thread information
available.
This function return true if it succeeded for every Thread, and false otherwise.
}

\begin{apient}
bool stopThreads() const
\end{apient}

\apidesc{
This function stops every Thread in this ThreadSet. It is similar to
Thread::stopThread.
This function return true if it succeeded for every Thread, and false
otherwise.
}

\begin{apient}
bool continueThreads() const
\end{apient}

\apidesc{
This function stops every Thread in this ThreadSet. It is similar to
Thread::continueThread.
This function return true if it succeeded for every Thread, and false
otherwise.
}

\begin{apient}
bool setSingleStepMode(bool v) const
\end{apient}

\apidesc{
This function puts every Thread in this ThreadSet into single step mode if v is
true. It clears single step mode if v is false. It is similar to
Thread::setSingleStepMode.
This function return true if it succeeded for every Thread, and false
otherwise.
}

\begin{apient}
bool getRegister(Dyninst::MachRegister reg,
                 std::map<Thread::ptr, Dyninst::MachRegisterVal> &res) const
\end{apient}

\apidesc{
This function gets the value of register reg in every Thread in this ThreadSet.
The collected values are returned in the res map, with each Thread mapped to
the value of reg in that thread. It is similar to
Thread::getRegister.
This function return true if it succeeded for every Thread, and false otherwise.
}

\begin{apient}
bool getRegister(Dyninst::MachRegister reg,
                 std::map<Dyninst::MachRegisterVal, ThreadSet::ptr> &res) const
\end{apient}

\apidesc{
This function gets the value of register reg in every Thread in this ThreadSet
and then aggregates all identical values together. The res map will contain
an entry for each unique register value, and map that value to a new ThreadSet
that contains every Thread that produced that register value. It is similar to
Thread::getRegister.
This function return true if it succeeded for every Thread, and false otherwise.
}

\begin{apient}
bool setRegister(Dyninst::MachRegister reg,
                 const std::map<ThreadSet::const_ptr, Dyninst::MachRegisterVal> &vals) const
\end{apient}

\apidesc{
This function sets the value of register reg in each Thread in this ThreadSet.
The value set in each thread is looked up in the vals map. It is similar to
Thread::setRegister.
This function\'s behavior is undefined if it is passed a
Thread that is not in this ThreadSet.
This function return true if it succeeded for every Thread, and false otherwise.
}

\begin{apient}
bool setRegister(Dyninst::MachRegister reg, Dyninst::MachRegisterVal val) const
\end{apient}

\apidesc{
This function sets the register reg to val in each Thread in this ThreadSet.
It is similar to Thread::setRegister.
This function return true if it succeeded for every Thread, and false otherwise.
}

\begin{apient}
bool getAllRegisters(std::map<Thread::ptr, RegisterPool> &results) const
\end{apient}

\apidesc{
This function gets the values of every register in each Thread in this
ThreadSet. The register values are returned as RegisterPools in the results
map, with each Thread mapped to its RegisterPool. It is similar to
Thread::getAllRegisters.
This function return true if it succeeded for every Thread, and false otherwise.
}

\begin{apient}
bool setAllRegisters(const std::map<Thread::const_ptr, RegisterPool> &val) const
\end{apient}

\apidesc{
This function sets the values of every register in each Thread in this
ThreadSet. The register values are extracted from the val map, with each
Thread specifying its register values via the map. This function is similar
to Thread::setAllRegisters.
This function\'s behavior is undefined if it is passed a
Thread that is not in this ThreadSet.
This function return true if it succeeded for every Thread, and false otherwise.
}

\begin{apient}
bool postIRPC(const std::multimap<Thread::const_ptr, IRPC::ptr> &rpcs) const
\end{apient}

\apidesc{
This function posts an IRPC to every Thread in this ThreadSet. The IRPC to
post to each Thread is specified in the rpcs multimap. This function is
similar to Thread::postIRPC.
This function return true if it succeeded for every Thread, and false otherwise.
}

\begin{apient}
bool postIRPC(IRPC::ptr irpc, std::multimap<Thread::ptr, IRPC::ptr> *result = NULL)
\end{apient}

\apidesc{
This function posts a copy of irpc to every Thread in this ThreadSet. If
result is non-NULL, then the new IRPC objects are returned in the result
multimap, with the Thread mapped to the IRPC that was posted there. This
function is similar to Thread::postIRPC.
This function return true if it succeeded for every Thread, and false otherwise.
}
